\section{The Real World: Historical Replay on NSE}
\label{sec:replay}

\subsection{Design}

The replay runs both dealers over the \RELIANCETestDays{} held-out days, minute by minute, with
the parameters from the calibration window. At the end of minute $t$ the dealer sees the close $S_t$,
its inventory $q_t$ and the time left, and posts a bid and an ask for minute $t+1$. The ask sells one
share at the quoted price if the next minute's high reaches it; the bid buys one share if the low
reaches it. This is the same fill rule used to calibrate $A$ and $k$, so the model is tested against
the fill process it was fitted to. The symmetric dealer again quotes AS's average half-spread for the
same day. Inventory left at 15:14 is marked at the last close. Nothing here is simulated: prices,
highs and lows are the real ones.

The dealer's PnL over a session splits exactly into three parts, minute by minute:
\begin{equation}
\Delta\text{PnL}_{t+1} =
\underbrace{\mathbb{1}^{\text{sell}}(a_t-S_t)+\mathbb{1}^{\text{buy}}(S_t-b_t)}_{\text{spread capture}}
+\underbrace{(\mathbb{1}^{\text{buy}}-\mathbb{1}^{\text{sell}})(S_{t+1}-S_t)}_{\text{adverse selection}}
+\underbrace{q_t\,(S_{t+1}-S_t)}_{\text{inventory carry}} ,
\label{eq:decomp}
\end{equation}
where $a_t,b_t$ are the quotes and the indicators mark fills in minute $t+1$. \emph{Spread capture}
is what the dealer earns relative to the mid at the moment it quotes. \emph{Adverse selection} is how
the mid moves in the same minute as a fresh fill: it is negative when the price keeps moving through
the dealer's quote after hitting it. \emph{Inventory carry} is the gain or loss on the position
already held. The model's assumption A6 says the second term is zero on average. The code checks
that the three parts add up to the PnL exactly.

\subsection{One Day Up Close}

Figure~\ref{fig:heroR} shows the test day on which the symmetric dealer's inventory strayed
furthest for RELIANCE. The stock rallied about 20 rupees in the first hour. The symmetric dealer
kept selling into the rally: it was short 14 shares by 11:00 and 23 at its extreme, and its PnL
fell to nearly $-400$ rupees at the worst point.
The AS dealer sold too, but each sale lowered both its quotes, so it bought back almost immediately
and its inventory never left $\pm3$. Panel (b) zooms into forty minutes to show the mechanism: the
blue band jumps up or down by a step every time a fill changes the inventory. Figure~\ref{fig:heroK}
shows the same for KSCL, where no-trade minutes leave flat stretches in the price.

\begin{figure}[!htbp]
\centering
\includegraphics[width=\textwidth]{replay_day_reliance}
\caption{RELIANCE, one held-out day. (a) Mid-price and every AS fill. (b) Zoom: AS ask and bid
(band), reservation price (dashed) and mid; triangles are fills at the quoted price. (c) Inventory
and (d) marked-to-market PnL of the two dealers.}
\label{fig:heroR}
\end{figure}

\begin{figure}[!htbp]
\centering
\includegraphics[width=\textwidth]{replay_day_kscl}
\caption{KSCL, one held-out day, in the same four views as Figure~\ref{fig:heroR}.}
\label{fig:heroK}
\end{figure}

\subsection{All Test Days}

Figure~\ref{fig:replayinv} overlays inventory on every test day. The picture is the same as in the
model's world (Figure~\ref{fig:mcpaths}): AS stays near zero, the symmetric dealer wanders up to
twenty shares away. Average absolute inventory is \RELIANCEReplayAbsqAS{} against
\RELIANCEReplayAbsqSym{} shares for RELIANCE and \KSCLReplayAbsqAS{} against \KSCLReplayAbsqSym{}
for KSCL.

\begin{figure}[!htbp]
\centering
\includegraphics[width=\textwidth]{replay_inventory_paths}
\caption{Inventory on every held-out day, real NSE minutes.}
\label{fig:replayinv}
\end{figure}

Risk control survives (Table~\ref{tab:replay}, Figure~\ref{fig:dailypnl}). The standard deviation
of daily PnL is \RELIANCEReplaySdAS{} rupees for AS against \RELIANCEReplaySdSym{} for the symmetric
dealer on RELIANCE, a factor of \RELIANCEReplaySdFactor{}, and \KSCLReplaySdAS{} against
\KSCLReplaySdSym{} on KSCL, a factor of \KSCLReplaySdFactor{}. Every AS day lands in a narrow range
around a small loss; the symmetric dealer's days scatter by hundreds of rupees.

\begin{figure}[!htbp]
\centering
\includegraphics[width=\textwidth]{replay_daily_pnl}
\caption{Daily PnL of the two dealers on each held-out day. The grey bar joins the two results for
the same day.}
\label{fig:dailypnl}
\end{figure}

\begin{table}[!htbp]
\centering
\caption{Historical replay, \RELIANCETestDays{} held-out days. Rupees per day (means), except the
standard deviation (S.d.) of daily PnL, inventory in shares and fills per day. Capture, Adverse and Carry
are the three parts of equation~\eqref{eq:decomp} and sum to PnL.}
\label{tab:replay}
\small
\setlength{\tabcolsep}{4.5pt}
\begin{tabular}{llrrrrrrrr}
\toprule
Stock & Strategy & PnL & S.d. & Capture & Adverse & Carry & Mean $|q|$ & Max $|q|$ & Fills \\
\midrule
RELIANCE & AS & -24.1 & 22.5 & +135.2 & -163.1 & +3.7 & 0.67 & 3.9 & 307 \\
RELIANCE & Symmetric & -78.0 & 108.8 & +151.6 & -183.5 & -46.2 & 6.69 & 15.0 & 249 \\
KSCL & AS & -28.9 & 19.4 & +80.1 & -128.2 & +19.3 & 0.48 & 3.4 & 135 \\
KSCL & Symmetric & +38.7 & 141.6 & +103.4 & -137.6 & +72.9 & 6.59 & 13.6 & 116 \\
\bottomrule
\end{tabular}

\end{table}

Profit does not survive. Both dealers lose money on average on RELIANCE (\RELIANCEReplayPnlAS{} and
\RELIANCEReplayPnlSym{} rupees a day), and Figure~\ref{fig:decomp} shows why. Spread capture is solidly
positive, \RELIANCECaptureAS{} rupees a day for AS, but adverse selection takes \RELIANCEAdverseAS{}:
on real NSE minutes, a quote is filled mostly when the price is moving through it, and the price
keeps going. Adverse selection is about the same size for both dealers, which makes sense: it is a
property of the order flow, not of where the quotes sit, and no amount of inventory skewing can
remove it. What skewing does remove is inventory carry: \RELIANCECarryAS{} rupees for AS, against
large swings for the symmetric dealer.

\begin{figure}[!htbp]
\centering
\includegraphics[width=\textwidth]{replay_pnl_decomposition}
\caption{PnL per day split into the three parts of equation~\eqref{eq:decomp}, mean $\pm$ one
standard error over the held-out days. Adverse selection, which the model assumes is zero, is the
largest term.}
\label{fig:decomp}
\end{figure}

On risk-adjusted terms the picture depends on the stock. For RELIANCE the certainty equivalent of the
daily PnL is \RELIANCEReplayCEAS{} for AS against \RELIANCEReplayCESym{} for the symmetric dealer. For KSCL it
is \KSCLReplayCEAS{} against \KSCLReplayCESym{}: the symmetric dealer comes out ahead. That result
rests on its inventory carry (+72.9 rupees a day, with a standard error of about 50), earned mostly
on two days (23 and 28 October) when the inventory it had drifted into happened to sit on the right
side of a large move. That is a lucky directional bet, not market making, and it is exactly the risk AS
is designed to avoid. Over all \KSCLDays{} days, including the calibration window, the symmetric
dealer's mean falls to \KSCLAllPnlSym{} rupees while its standard deviation stays about
\KSCLAllSdFactor{} times AS's. The RELIANCE results are the same over all days (AS \RELIANCEAllPnlAS{},
symmetric \RELIANCEAllPnlSym{}, standard deviation ratio \RELIANCEAllSdFactor{}).
